\section{Verdict: Good Project, or National Need?}\label{sec:verdict}

\begin{finding}{The evidence supports a two-stage answer.}
\textbf{Project quality is the entry ticket; national timing is the selector.}
Nothing in the record shows a weak project winning: every examined winner had
a working artifact, used open data, and could demonstrate. But nothing in the
record shows raw technical excellence winning \emph{against} the year's
current either --- and the one winner whose technical depth most clearly
exceeded its peers (A/B Street) won in precisely the year its subject matter
became the zeitgeist, three years after it was built. The rubric's arithmetic
(70\% of preliminary weight on feasibility and social importance, 30\% on
innovation\cite{handbook}), the annually re-set presidential theme, and the
diplomatic staging of the award all point the same way as the case record.
\end{finding}

Put more bluntly, for the benefit of our own entry: the question ``is it
really a good project or just the nation's immediate need?'' assumes a
dichotomy the competition does not contain. The Presidential Hackathon is a
\emph{policy showcase that uses competition mechanics}. It selects for the
intersection: projects good enough to demo without embarrassment, aimed
precisely at what the state most needs demonstrated that year, carried by
teams whose composition serves Taiwan's international story. Within that
intersection, ``better project'' still matters --- the 2025 field had nine
finalists\cite{focuscrop}, and demonstrable maturity visibly separates winners
from the invisible losing pool --- but outside that intersection, quality
alone has no observed path to victory.

Two honest caveats bound this verdict. First, survivorship: we cannot examine
rejected entries, so ``no weak project wins'' is really ``no weak project is
visible among winners.'' Second, afterlife is a lagging and noisy indicator:
CropNow and Beyond Hearing are recent, and quiet teams may be active teams.
The 50/50 afterlife split (Finding~3) should be read as ``the competition does
not guarantee sustainability,'' not as ``half of all winners failed.''

\subsection{What this means for the CarbonPass entry}

The findings translate directly into entry strategy, and they largely confirm
--- with one sharpened warning --- the approach already taken in this
project's proposal documents:

\begin{itemize}[leftmargin=1.3em]
\item \textbf{Confirmed: theme-lock is non-negotiable.} No winner has beaten
the year's current. CarbonPass's framing of the carbon-data divide as 2026's
digital-inclusion problem is the correct kind of move --- provided the
inclusion is visible in the demo (the LINE/voice interface), not just argued
in prose (Finding~1).
\item \textbf{Confirmed: bring something that already works.} Half the
winners pre-existed the event; feasibility outweighs innovation in the rubric.
The plan to arrive with a working data-pack engine on real bill formats, and
to secure a Gangshan pilot during mentorship, mirrors what distinguishes
winners (Finding~2).
\item \textbf{Confirmed: the multinational team and the exportable story
matter.} India's recurrence and the ceremony's diplomatic cast are data
(Finding~4). CarbonPass's ``every SME-exporting economy faces this wall''
expansion path is the right register.
\item \textbf{Sharpened warning: winning is not a sustainability plan.} The
afterlife record says the competition amplifies but does not sustain
(Finding~3). Every winner that lasted had an institution behind it. For
CarbonPass, the integration targets (MOENV's CBAM platform, SMESA's service
network, TIFI/MIRDC) are therefore not a rhetorical flourish in the proposal
--- they are, on this evidence, the difference between becoming CoST Honduras
and becoming Re-Fill City.
\end{itemize}

\section{Limitations}\label{sec:limits}

This review covers the International Track's publicly celebrated teams (two
per year in official English-language reporting); some editions named
additional excellent teams not consistently listed in English. Afterlife
searches were conducted in English in July 2026 and may miss activity
documented only in Mandarin, Thai, or Bahasa. The 2019 honor roll contains a
source discrepancy we flag but cannot resolve (Section~\ref{sec:corrections}).
Assessments of ``technical substance'' are qualitative readings of documented
functionality, not code review.
